
\subsubsection{6.1 Interpretation of Results}

\textbf{Attention concentration as a structural property within domain.} The 100\% satisfaction rate (Gini > 0.5 and top-10\% share > 0.5 across all 200 evaluation examples) indicates that heavy-hitter attention concentration is a consistent property of Llama-2-7B's trained weights within the WikiText-103 evaluation domain. The low variance (Gini std = 0.0117) suggests that this is not driven by specific inputs but is a relatively stable feature of the model's attention distributions on this domain's text. Cross-domain generalization --- whether the same concentration pattern appears on code, scientific text, or dialogue --- remains untested.

\textbf{The observed PPL change in h-e2 ($\Delta = -0.016)$.} The entropy-guided k = 4 SWA conversion produced a negligible PPL decrease rather than an increase. A decrease of this magnitude (0.016 on a baseline of 6.34) is within expected measurement noise for strided perplexity evaluation and should not be interpreted as the SWA conversion improving the model. The result is more precisely characterized as: converting the 4 highest-entropy layers to SWA(w = 512) causes no measurable perplexity degradation on WikiText-103. This is consistent with the hypothesis that high-entropy layers are not exploiting long-range global attention in ways that are critical for WikiText-103 language modeling.

The selected layers (0, 1, 10, 31) include both the first two layers of the model (which in transformer language models typically handle low-level token and positional features rather than deep semantic integration) and layers at mid and late depth. The concentration of the highest-entropy scores in layers 0 and 1 (3.201 and 2.972, substantially above all others) suggests that early-depth layers in Llama-2-7B have particularly diffuse attention patterns, which is consistent with prior findings about the role of early transformer layers.

\textbf{Head-mean pooling as a methodological finding.} The 32\% gap between head-mean Gini (0.681) and head-max Gini (0.466) has consequences beyond this paper. For any entropy-based method that characterizes transformer layers --- whether for pruning, routing, or structural modification --- the choice of head aggregation method is not an implementation detail. Head-max is dominated by outlier heads and underestimates the concentration shared across a layer's heads. Head-mean captures the layer-level structural property. This distinction should be explicitly justified in future work using per-layer entropy signals.

\textbf{Layer selection depth distribution.} The top-4 entropy-selected layers include 2 early-depth layers (0, 1), 1 mid-depth layer (10), and 1 late-depth layer (31). The large gap between layers 0--1 and all other layers (entropy $\approx$ 3.2 and 2.97 vs. next highest $\approx$ 1.69) suggests that the early layers of Llama-2-7B are qualitatively different from deeper layers in their attention spread. Whether this depth pattern generalizes to k = 8 conversion or holds across calibration domains is an open question.

\subsubsection{6.2 Open Questions from Pending Experiments}

\textbf{Does entropy-guided selection outperform uninformed selection?} The h-e2 result confirms that entropy-guided k = 4 SWA is quality-preserving, but it does not establish that entropy-guided selection is superior to random or last-k selection. The preliminary QA F1 probe (Section 5.4) provides a directional but non-significant signal (0.43 pp vs. 0.67 pp degradation, p = 0.4507). The proper comparison requires h-m1: identical SWA conditions (w = 512, k = 4) with random and last-k layer selection, measuring WikiText-103 perplexity at adequate statistical power.

\textbf{How does perplexity degrade at k = 8?} Whether converting 8 of 32 layers (25\%) causes measurable degradation beyond the 2-point threshold is untested. The k = 8 case would provide the first characterization of the feasibility boundary for zero-shot selective SWA in Llama-2-7B. Given that layers 0 and 1 dominate the entropy ranking, k = 8 would extend selection into layers with substantially lower entropy scores (< 1.7), where SWA compatibility may differ.

\textbf{What is the residual stream compensation mechanism?} The theoretical account for why SWA layers at selected positions do not degrade quality is that the remaining 28 full-attention layers compensate via the residual stream. This mechanism is not directly tested. Depth-position analysis in h-m2 could provide indirect evidence: if early-depth SWA layers degrade quality more than late-depth ones, it would suggest that residual stream buildup matters for compensation.

\subsubsection{6.3 Limitations}

\textbf{L1: Entropy selection superiority over baselines is unconfirmed.} The h-m1 comparison (entropy vs. random vs. last-k) has not been executed. The h-e2 result confirms that entropy-guided selection is quality-preserving; it does not confirm that it is necessary for quality preservation --- random k = 4 SWA might also be within the 2-point threshold. This is the central remaining open question.

\textbf{L2: Single model and single evaluation domain.} All results apply to Llama-2-7B evaluated on WikiText-103. Generalizability to other model architectures (Llama-3 with grouped-query attention, encoder-decoder models), other scales, or other evaluation domains is untested. The entropy concentration finding may be domain-specific.

\textbf{L3: Proxy QA F1 evidence (Section 5.4) is methodologically misaligned.} The exploratory QA probe uses top-k attention score retention (not SWA masking) and F1 (not perplexity). The directional trend (entropy better: 0.43 pp vs. 0.67 pp) is reported for transparency but cannot support claims about entropy-guided SWA superiority.

\textbf{L4: Exact per-pair Spearman $\rho$ values are not tabulated.} The gate criterion (min $\rho \geq$ 0.8) is confirmed, but exact point estimates per subset pair are not available in the h-e1 output files. Reviewers cannot assess whether stability is narrowly above threshold or substantially above it. These values should be extracted from h-e1 implementation logs before final submission.

\textbf{L5: Mechanism verification ambiguity.} The first h-e2 run produced a mechanism verification warning (''Layer 0: expected 512 attended positions, got 600''), interpreted in the second run as measuring the pre-patch mask state via a forward hook. The second run verified mechanism activation via PPL delta (|$\Delta |$ = 0.016 $\neq$ 0), and swa\_patch.py confirms correct mask construction semantics (col <= row and col >= row $-$ window\_size + 1). Nonetheless, the mechanism verification ambiguity represents an implementation risk that warrants explicit logging of attended positions in any follow-up experiment.

\textbf{L6: Calibration terminated early.} The h-e2 calibration completed after 63 of 100 requested sequences. The entropy ranking [0, 1, 31, 10] was produced from this partial calibration. Whether completion of the full 100-sequence calibration would change the top-4 selection is not verified; the ranking stability analysis (h-e1) suggests robustness, but this specific early termination was not part of the designed stability validation.

\subsubsection{6.4 Broader Impact}

This work targets inference efficiency for large language models without fine-tuning. The entropy scoring requires only 100 calibration sequences and a single GPU, making it accessible to practitioners with limited compute. If the selective SWA conversion generalizes beyond Llama-2-7B, it could reduce inference cost for models deployed at long contexts, potentially lowering energy consumption and enabling deployment on hardware with less GPU memory. Practitioners should validate that SWA conversion does not disproportionately affect performance on specific tasks, languages, or domains before deployment in production settings, as the current validation is limited to WikiText-103 perplexity.

